\section{Where do binning and FAST struggle?}

Scalar action binning gives each action dimension at each timestep its own
discrete symbol. It is simple, but chunking high-frequency actions can create
long sequences. OpenVLA is a concrete example of per-dimension discretization;
its tokenizer implementation is linked in the references.

FAST addresses redundancy through a frequency-based representation using the
discrete cosine transform, followed by quantization and token compression.
Importantly, the FAST paper reports strong results, including settings where
scalar discretization struggles. Its success supports the idea that
representation matters.

Our question concerns a possible remaining limitation in the single-transformer
setting: does the resulting token sequence concentrate predictive difficulty
into a few positions? The draft proposes that early tokens may carry difficult
decisions while later tokens contribute less useful information. This is a
hypothesis to measure, not an established consequence of the DCT. Compression,
ordering, data, and conditioning all affect the final token distribution.

The relevant quantity is conditional uncertainty: how hard is a token to predict
given the observation, instruction, and preceding tokens? A broad token
histogram alone does not answer that question, and high-frequency motion is not
necessarily noise.

Exposure bias adds a second issue. During teacher-forced training, each token is
predicted after an expert prefix. During generation, it follows the model’s own
predictions. An early mistake can change the context for later decisions,
although one wrong token need not invalidate an entire action chunk.

A useful comparison measures sequence length, loss by token position,
reconstruction error, and closed-loop task success together. None of these
measurements can replace the others.
